\documentclass[11pt]{article}
\usepackage[margin=1in]{geometry}
\usepackage{amsmath,amssymb,amsthm}
\usepackage{physics}
\usepackage{enumitem}
\usepackage{hyperref}
\usepackage{titlesec}

\theoremstyle{definition}
\newtheorem{definition}{Definition}[section]
\newtheorem{example}{Example}[section]
\newtheorem{remark}{Remark}[section]
\newtheorem{proposition}{Proposition}[section]
\newtheorem{theorem}{Theorem}[section]

\titleformat{\section}{\large\bfseries}{Lecture 13.\thesection}{1em}{}

\title{\textbf{Lecture 13: Correcting Quantum Errors} \\ \large Understanding Quantum Information \& Computation}
\author{Based on the lectures by John Watrous (IBM Quantum)}
\date{}

\begin{document}
\maketitle
\tableofcontents
\newpage

%----------------------------------------------------------------------
\section{Overview}
%----------------------------------------------------------------------

This lecture opens the fourth and final unit of the series: \textbf{Quantum Error Correction}.  The central question is:

\begin{center}
\emph{Can we protect quantum information against noise, even though we cannot copy quantum states?}
\end{center}

The answer is \emph{yes}, and the key idea is to encode quantum information \emph{redundantly} into entangled states of multiple qubits.  Topics:

\begin{enumerate}
    \item Classical repetition codes (warm-up).
    \item Why na\"ive repetition fails for qubits.
    \item Quantum error-correcting codes: the 3-qubit bit-flip code and the 3-qubit phase-flip code.
    \item The Shor (9-qubit) code as the first code correcting arbitrary single-qubit errors.
\end{enumerate}

%----------------------------------------------------------------------
\section{Classical Repetition Codes}
%----------------------------------------------------------------------

\subsection{The Idea}

Protect a classical bit against a single bit-flip error by encoding it redundantly:
\[
    0 \;\longrightarrow\; 000, \qquad 1\;\longrightarrow\; 111.
\]
If at most one of the three bits is flipped, \textbf{majority vote} recovers the original bit.

\subsection{Formalisation}

\begin{definition}[Repetition code]
The \textbf{$n$-bit repetition code} encodes one logical bit into $n$ physical bits:
\[
    0\mapsto \underbrace{00\cdots 0}_{n},\qquad 1\mapsto \underbrace{11\cdots 1}_{n}.
\]
It can correct up to $\lfloor(n-1)/2\rfloor$ bit-flip errors via majority decoding.
\end{definition}

%----------------------------------------------------------------------
\section{Obstacles to Quantum Error Correction}
%----------------------------------------------------------------------

Three features of quantum information make error correction seem impossible at first:
\begin{enumerate}
    \item \textbf{No-cloning theorem:} we cannot copy an unknown quantum state, so na\"ive repetition $\ket{\psi}\to\ket{\psi}\ket{\psi}\ket{\psi}$ is forbidden.
    \item \textbf{Continuous errors:} quantum errors are not just bit flips; a qubit can be rotated by an arbitrary (continuous) amount.
    \item \textbf{Measurement disturbance:} measuring a quantum state to diagnose errors generally destroys the encoded information.
\end{enumerate}

All three obstacles can be overcome.

%----------------------------------------------------------------------
\section{The 3-Qubit Bit-Flip Code}
%----------------------------------------------------------------------

\subsection{Encoding}

We encode \emph{one logical qubit} into \emph{three physical qubits}:
\[
    \ket{0}_L = \ket{000}, \qquad \ket{1}_L = \ket{111}.
\]
An arbitrary logical state $\alpha\ket{0}+\beta\ket{1}$ is encoded as
\[
    \ket{\psi}_L = \alpha\ket{000}+\beta\ket{111}.
\]
Note: this is \emph{not} cloning $\ket{\psi}$.  The state $\alpha\ket{000}+\beta\ket{111}$ is entangled and differs from $(\alpha\ket{0}+\beta\ket{1})^{\otimes 3}$.

The encoding circuit uses two CNOT gates: the logical qubit controls two ancilla qubits initialised to $\ket{0}$.

\subsection{Error Model}

We assume at most one of the three qubits suffers a bit-flip ($X$ error).  There are four cases: no error, $X$ on qubit~1, $X$ on qubit~2, $X$ on qubit~3.

\subsection{Syndrome Measurement}

\begin{definition}[Syndrome]
The \textbf{syndrome} is classical information extracted by measuring certain operators that reveals \emph{which} error occurred, without revealing any information about the encoded state $\alpha,\beta$.
\end{definition}

For the 3-qubit bit-flip code, the syndrome operators are $Z_1 Z_2$ and $Z_2 Z_3$ (parity checks on pairs of qubits).

\begin{center}
\begin{tabular}{ccc}
\hline
\textbf{Error} & $Z_1 Z_2$ & $Z_2 Z_3$ \\\hline
None ($I$) & $+1$ & $+1$ \\
$X_1$ & $-1$ & $+1$ \\
$X_2$ & $-1$ & $-1$ \\
$X_3$ & $+1$ & $-1$ \\
\hline
\end{tabular}
\end{center}

Each error produces a unique syndrome $\to$ the error can be identified and corrected by applying the appropriate $X$ gate.

\paragraph{Key point.}  The syndrome measurements are \emph{joint} measurements on pairs of qubits.  They reveal parity information but not the values of $\alpha$ and $\beta$, so the encoded state is preserved.

\subsection{Limitations}

The 3-qubit bit-flip code corrects only $X$ errors.  It cannot handle phase-flip ($Z$) errors.

%----------------------------------------------------------------------
\section{The 3-Qubit Phase-Flip Code}
%----------------------------------------------------------------------

\subsection{Encoding}

\[
    \ket{0}_L = \ket{+++}, \qquad \ket{1}_L = \ket{---},
\]
or equivalently $\ket{0}_L = H^{\otimes 3}\ket{000}$, $\ket{1}_L = H^{\otimes 3}\ket{111}$.

\subsection{Error Correction}

A $Z$ error in the $\{\ket{+},\ket{-}\}$ basis acts like a bit flip: $Z\ket{+}=\ket{-}$, $Z\ket{-}=\ket{+}$.

Syndrome operators: $X_1 X_2$ and $X_2 X_3$ (parity checks in the $X$-basis).

\subsection{Limitations}

Corrects only $Z$ errors.  Cannot handle $X$ errors.

%----------------------------------------------------------------------
\section{The Shor Code (9-Qubit Code)}
%----------------------------------------------------------------------

\subsection{Construction}

Shor's code \emph{concatenates} the phase-flip code with the bit-flip code:
\begin{enumerate}
    \item First, encode against phase flips: $\ket{0}_L\to\ket{+++}$, $\ket{1}_L\to\ket{---}$.
    \item Then, encode each $\ket{+}$ and $\ket{-}$ against bit flips using the 3-qubit repetition: $\ket{+}\to\frac{1}{\sqrt{2}}(\ket{000}+\ket{111})$, $\ket{-}\to\frac{1}{\sqrt{2}}(\ket{000}-\ket{111})$.
\end{enumerate}

Result: 1 logical qubit encoded into 9 physical qubits.
\begin{align}
    \ket{0}_L &= \frac{1}{2\sqrt{2}}\bigl(\ket{000}+\ket{111}\bigr)^{\otimes 3}, \\
    \ket{1}_L &= \frac{1}{2\sqrt{2}}\bigl(\ket{000}-\ket{111}\bigr)^{\otimes 3}.
\end{align}

\subsection{Error Correction Capability}

\begin{theorem}
The Shor code can correct any single-qubit error on any one of the 9 physical qubits.  This includes:
\begin{itemize}
    \item Bit-flip errors ($X$),
    \item Phase-flip errors ($Z$),
    \item Combined bit-and-phase errors ($Y=iXZ$),
    \item Arbitrary single-qubit errors (by linearity).
\end{itemize}
\end{theorem}

\paragraph{Handling continuous errors.}
Any single-qubit error can be written $E = \alpha_0 I + \alpha_1 X + \alpha_2 Y + \alpha_3 Z$.  The syndrome measurement \emph{projects} the error onto one of the discrete Pauli errors $\{I,X,Y,Z\}$.  After syndrome measurement and correction, the encoded state is restored exactly.  This is why continuous errors do not pose a fundamental obstacle.

%----------------------------------------------------------------------
\section{The Error Discretisation Principle}
%----------------------------------------------------------------------

\begin{theorem}[Error discretisation]
If a quantum error-correcting code can correct a set of errors $\{E_1,\dots,E_r\}$, then it can correct \emph{any} error in the linear span of $\{E_1,\dots,E_r\}$.
\end{theorem}

Since $\{I,X,Y,Z\}$ spans all $2\times 2$ matrices, a code that corrects all single-qubit Pauli errors automatically corrects \emph{all} single-qubit errors (including continuous rotations).

%----------------------------------------------------------------------
\section{Summary}
%----------------------------------------------------------------------

\begin{itemize}
    \item Quantum error correction is possible despite no-cloning, continuous errors, and measurement disturbance.
    \item The \textbf{3-qubit bit-flip code} corrects $X$ errors; the \textbf{3-qubit phase-flip code} corrects $Z$ errors.
    \item The \textbf{Shor (9-qubit) code} concatenates the two and corrects \emph{any} single-qubit error.
    \item \textbf{Syndrome measurements} diagnose errors without disturbing the encoded information.
    \item The \textbf{error discretisation principle} reduces continuous errors to discrete Pauli errors.
\end{itemize}

\end{document}
